% !TEX root = ../mainfile.tex
% !TEX encoding = UTF-8 Unicode
% !TEX spellcheck = en_US

\chapter{Background}
\label{ch:background}

This chapter documents the episodes the empirical chapters analyse and establishes the dates that define the event windows of Chapter~\ref{ch:methodology}. For each episode it sets out what happened and how official institutions described the market reaction.

\section{Liberation Day}
\label{sec:bg_tariff}

On 2~April 2025, in an announcement branded ``Liberation Day,'' the United States imposed broad ``reciprocal'' tariffs on nearly all of its trading partners, combining a flat baseline rate of ten percent with higher country-specific surcharges. The measure was imposed under the International Emergency Economic Powers Act, after the President declared that large and persistent trade deficits and a lack of reciprocity in trade relationships amounted to a national emergency. It applied to most U.S.\ imports of goods \parencite{crs_2025_tariff_timeline}. Later assessments found that foreign exporters absorbed only about five percent of the tariff cost, so almost all of it fell on U.S.\ firms and consumers, while import volumes declined sharply \parencite{ecb_2026_bulletin}.

Equity prices fell sharply and implied volatility spiked to a level exceeded only during the global financial crisis and the COVID-19 pandemic \parencite{ecb_2025_fsr}, while crude oil fell as markets priced slower growth and weaker oil demand \parencite{eia_2026_crude}. Around the announcements the dollar depreciated against the other major advanced-economy currencies even as risk metrics deteriorated and Treasury yields rose \parencite{kamin_2025}. A partial ninety-day pause followed on 9~April 2025, suspending the country-specific surcharges while the baseline rate stayed in force \parencite{crs_2025_tariff_timeline}.

The legal basis of the tariffs remained contested. Multiple federal courts ruled that the President had exceeded his authority in at least some of the actions. The tariffs stayed in force while the Supreme Court reviewed the cases \parencite{crs_2025_tariff_timeline}. On 20~February 2026 the Supreme Court held that the act does not authorise the President to impose tariffs \parencite{scotus_2026_tariffs}. Tariffs imposed under other trade statutes were not before the Court.

\section{Twelve-Day War and Hormuz Crisis}
\label{sec:bg_iran}

Israel struck Iran on 13~June 2025, opening a twelve-day war that drew in U.S.\ strikes on three Iranian nuclear sites on 22~June. It ended on 24~June, hours after an Iranian missile strike on a U.S.\ base in Qatar \parencite{crs_2025_israel_iran}. The episode raised the risk of a Gulf supply disruption and lifted crude only briefly \parencite{kilian_plante_richter_zhou_2026,eia_2026_crude}.

The conflict resumed on 28~February 2026, when the United States and Israel opened a joint campaign of airstrikes on Iran \parencite{crs_2026_iran_conflict}. Both governments said they aimed to end Iran's nuclear and missile capability and to force a change of government in Tehran. Iran's response was regional, reaching Israel, U.S.\ military sites across the Middle East, and the Arab states that host American forces, among them Bahrain, Kuwait and Qatar. Iran also announced that it would close the Strait of Hormuz \parencite{hoc_2026_conflict}, which carries roughly a quarter of globally traded oil \parencite{fattouh_economou_2026}. Attacks on shipping and the repricing of tanker insurance effectively closed the Strait within days \parencite{kilian_plante_richter_zhou_2026}. On 13~March 2026 the United States reportedly struck military targets on Kharg Island, the terminal through which some ninety percent of Iran's oil exports transit \parencite{crs_2026_iran_conflict}. Brent crude, which had traded near \$71 per barrel before the 28~February strikes, rose by 46~percent over March to a monthly average of about \$104 and, in front-month futures, peaked at \$118 on 29~April \parencite{fattouh_economou_2026,imf_2026_menap,eia_2026_hormuz}. The daily spot series plotted in Figure~\ref{fig:ov_oil} reached \$138.21 on 7~April \parencite{ds_eia_energy}.

A two-week ceasefire announced on 7~April 2026 came with an Iranian declaration that the Strait was open, though only to vessels using a route coordinated with Iranian port authorities, which falls short of free passage \parencite{fattouh_economou_2026}. Talks in Islamabad on 11 and 12~April broke down, and the United States began a naval blockade of Iranian ports on 13~April. The ceasefire itself held into May, but Iran resumed attacks on 4~May and the United States struck Iran again on 25~May \parencite{crs_2026_ceasefire,cnn_2026_strikes}. A memorandum of understanding signed on 17~June set out plans for talks over an initial sixty-day period on freedom of navigation, Iran's nuclear and missile programmes and sanctions \parencite{hoc_2026_ceasefire}. Once it was signed and tanker movements picked up, Brent fell through the rest of the quarter to \$72 on 26~June \parencite{eia_2026_hormuz}.

The fighting did not stop with the agreement. Between 25 and 28~June the United States struck Iran in response to an attack on commercial shipping in the Strait, and Iran struck U.S.\ bases in Bahrain and Kuwait \parencite{hoc_2026_ceasefire}. The sample analysed in this thesis ends on 30~June 2026, with the conflict unresolved. Israel's part in the fighting was concentrated in the first weeks. An exchange of fire with Iran on 8~June was its first since the April ceasefire, and both sides then backed away from further strikes \parencite{ap_2026_israel_iran}. Israel was a partner in the military campaign but had no part in shaping the memorandum of 17~June \parencite{inss_2026_mou}, and it sat out the fighting between the United States and Iran in July \parencite{toi_2026_sidelines}.

\section{Russian Invasion of Ukraine}
\label{sec:bg_ukraine}

On 24~February 2022 Russian forces crossed the border into Ukraine and began a full-scale assault on the country. The European Union, the United Kingdom, the United States and other allies answered within days with coordinated sanctions targeting Russia's banks, its aviation and shipping sectors, and industries of strategic significance to the Russian state, among them defence, aerospace and energy \parencite{hoc_2025_sanctions}.

Two measures reached the international payments and reserve system directly. On 26~February 2022 the European Union, the United Kingdom, the United States and Canada announced that they would restrict the Russian Central Bank's access to its foreign reserves. The stated aim was to undercut its ability to conduct foreign exchange transactions in support of the rouble. Russia was thought to hold more than \$630~billion in currency and gold reserves at the time. By late 2023 more than sixty percent of them had been frozen. On 2~March 2022 the European Union announced the exclusion of seven Russian banks from the SWIFT (Society for Worldwide Interbank Financial Telecommunication) payments network, effective ten days later. Three more followed in June \parencite{hoc_2025_sanctions}.

Energy markets had already been pricing the risk of an invasion. Crude oil rose and grew more volatile from mid-January 2022 onward, passed \$100 per barrel as the attack began, and Brent futures closed just below \$115 per barrel on 2~March \parencite{eia_2022_ukraine}.

\section{Overview of the Events}
\label{sec:bg_design}

The analysis dates six events. The two treatments are Liberation Day, a trade-policy shock the United States authored, on 2~April 2025, and the Hormuz crisis, a military and oil-supply shock in which it was a belligerent, on 28~February 2026. Each is followed by a reversal within weeks, the ninety-day tariff pause on 9~April 2025 and the two-week ceasefire on 7~April 2026. The invasion of Ukraine on 24~February 2022 enters as the external control, a shock the United States did not author, and the twelve-day war of 13~June 2025 enters as supporting context, reported alongside the Hormuz crisis but outside the main formal cross-event comparison.

Figures~\ref{fig:ov_dollar}--\ref{fig:ov_gpr} place these events in their longer-run context, plotting the broad dollar index, the price of Brent crude, and the geopolitical-risk index of \textcite{caldara_iacoviello_2022} over 2019--2026, with the four main events marked on each series. Chapter~\ref{ch:results} measures the responses formally.

\clearpage

\begin{figure}[p]
  \centering
  \includegraphics[width=\linewidth]{fig_overview_dollar}
  \caption[Broad dollar index, 2019--2026]{Broad dollar index, 2019--2026. Source: \textcite{ds_fed_dollar}, via FRED (Federal Reserve Economic Data).}
  \label{fig:ov_dollar}

  \vspace{50pt}
  % \vspace{60pt}   % superseded 2026-09-09 (float overflowed by 8.8pt)
  \includegraphics[width=\linewidth]{fig_overview_oil}
  \caption[Brent crude oil price, 2019--2026]{Brent crude oil price, 2019--2026, USD per barrel. Source: \textcite{ds_eia_energy}, via FRED.}
  \label{fig:ov_oil}

  \vspace{50pt}
  % \vspace{60pt}   % superseded 2026-09-09 (float overflowed by 8.8pt)
  \includegraphics[width=\linewidth]{fig_overview_gpr}
  \caption[Geopolitical Risk Index, 2019--2026]{Daily Geopolitical Risk Index, 2019--2026. Source: \textcite{ds_gpr}.}
  \label{fig:ov_gpr}
\end{figure}
